\chapter{偏微分方程基本概念}
\intro{本章介绍偏微分方程的分类、定解条件和适定性理论。}

\begin{mydefinition}{偏微分方程}
含有未知函数 $u(x_1,\ldots,x_n)$ 及其偏导数的方程。最高阶偏导数的阶数称为方程的\textbf{阶}。
\end{mydefinition}

\begin{mydefinition}{二阶线性偏微分方程}
\[
\sum_{i,j=1}^{n}a_{ij}\frac{\partial^2 u}{\partial x_i\partial x_j}+\sum_{i=1}^{n}b_i\frac{\partial u}{\partial x_i}+cu=f
\]
\end{mydefinition}

\section{方程的分类}

\begin{mydefinition}{特征值判别}
对于 $Au_{xx}+2Bu_{xy}+Cu_{yy}+\cdots=0$，判别式 $\Delta=B^2-AC$：
\begin{itemize}
	\item $\Delta>0$：双曲型（波动方程）
	\item $\Delta=0$：抛物型（热传导方程）
	\item $\Delta<0$：椭圆型（拉普拉斯方程）
\end{itemize}
\end{mydefinition}

\section{定解条件}

\begin{mydefinition}{初始条件}
\begin{itemize}
	\item 热传导方程（一阶时间导数）：$u(x,0)=\phi(x)$
	\item 波动方程（二阶时间导数）：$u(x,0)=\phi(x)$，$u_t(x,0)=\psi(x)$
\end{itemize}
\end{mydefinition}

\begin{mydefinition}{三类边界条件}
\begin{enumerate}
	\item Dirichlet：$u|_{\partial\Omega}=g$
	\item Neumann：$\displaystyle\frac{\partial u}{\partial n}\bigg|_{\partial\Omega}=g$
	\item Robin：$\displaystyle\left(\alpha u+\beta\frac{\partial u}{\partial n}\right)\bigg|_{\partial\Omega}=g$
\end{enumerate}
\end{mydefinition}

\section{适定性与叠加原理}

\begin{mydefinition}{适定性（Well-posedness）}
解\textbf{存在}、解\textbf{唯一}、解\textbf{稳定}（连续依赖于初边值条件）。
\end{mydefinition}

\begin{theorem}[叠加原理]
若 $u_1,\ldots,u_n$ 是线性齐次方程的解，则 $u=\sum_{i=1}^{n}c_i u_i$ 也是解。
\end{theorem}

\begin{remark}
叠加原理是分离变量法和格林函数法的理论基础。
\end{remark}
